\chapter{Reliability and manufacturing at scale}
\label{ch:reliability}

\textit{Read first:} FIT versus DPPM; qualification argument; sample confidence;
ATP versus qualification versus SPC; NPI; escaped defects.

\textit{Deep dive:} Arrhenius life projection; GR-468 stress detail; connector
mate-cycle physics.

\textit{Reference:} wear-out map, NPI gate table, ATP control checklist,
yield-owner split.

A link that closes in the lab can still fail the business case if lasers die in
the field or suppliers cannot hold yield. At gigawatt, multi-generation scale,
reliability and manufacturability stop being afterthoughts and become design
constraints: they decide whether you put the laser on the ASIC package or in a
replaceable module, how hard you derate, and what ATP language you freeze with
partners. This chapter covers the vocabulary of failure at scale, the
qualification flows that project field life, and the supplier-execution work these
systems demand.

\section{The language of scale: FIT and DPPM}
\label{sec:fit-dppm}

Fleet arguments need two different numbers. One is about life in the field; the
other is about quality at the factory door. A rate without population
definition, observation time, and confidence is not a fleet claim
(\Cref{sec:tree-evidence-block}).

\begin{description}
\item[FIT (failures in time)]\firstterm{FIT} failures per $10^{9}$ device-hours.
      State the population unit (laser die, module, or link), observed versus
      predicted FIT, confidence interval, useful-life / constant-hazard
      assumption, and whether the system is treated as repairable. Multiply a
      per-unit FIT by population and hours only after those definitions are
      fixed.
\item[DPPM (defective parts per million)] Distinguish incoming defects, outgoing
      defects, and escaped field defects. Always state the denominator, time
      window, and whether rework and retest are included.
\end{description}

\keyidea{Zero failures do not prove a zero failure rate.
0 failures in 20 units is not the same evidence as 0 failures in 1{,}000 units.
0 failures after $10^{3}$ device-hours is not the same evidence as 0 failures
after $10^{6}$ device-hours. Evidence strength depends on sample-hours,
one-sided upper confidence bounds, censoring, and whether lots and sites are
representative. Qualification and production SPC complement each other; neither
alone is a fleet claim.}

\aside{An order-of-magnitude exercise: a
$100{,}000$-GPU cluster with several optical links per GPU has on the order of
$10^{5}$--$10^{6}$ lasers. Even a small per-laser FIT then implies a nontrivial
number of link failures per day, which is exactly why field-replaceable external
laser sources and redundancy are attractive.}

\section{Qualification flows}

Qualification is a life and variation argument, not a list of stresses. A
requirement such as five-year operation is too abstract to test directly. You
identify the mechanisms that could violate it, choose stresses that accelerate
those mechanisms, monitor an observable signature, and decide whether the
resulting evidence is sufficient for the claim
(\Cref{sec:tree-qual-evidence,sec:interview-validation-ladder}).

Debugging is what you do when remaining margin hits zero. Qualification asks
how much margin remains after the expected stresses. Use the same canonical
validation lifecycle as \Cref{sec:tree-qualification}; this chapter owns the
environmental, reliability, manufacturing, and ATP gates. Keep the customer view
and the vendor view distinct: the vendor designs internals; the customer
characterizes externally visible behavior and decides deployment
(\Cref{sec:interview-margin-budget,sec:interview-customer-vendor}).

\paragraph{One continuous example: gradual laser degradation.}
Suppose the concern is gradual laser degradation. HTOL is useful only because
elevated temperature and operation may accelerate the active-region mechanism.
LIV measurements before and during stress reveal whether threshold current rises
or slope efficiency falls. Sample size, lots, stress hours, and confidence
determine the strength of the resulting life claim
(\Cref{sec:qual-sample-strategy,sec:fit-dppm}). A production burn-in may later
screen infant mortality, but it does not replace qualification.

\begin{dectree}
Requirement
  |
Budget (life / FIT / DPPM / power)
  |
Allocation (die / package / module / host)
  |
Verification (GR-468 / GR-1221 / JESD47 / ATP)
  |
Production + field monitoring
\end{dectree}

\engcheck{Margin budgeting}{%
Every stress spends margin: temperature, voltage, ripple, contamination,
insertion loss, connector wear, vibration, aging, process variation.
Qual verifies what remains, not only that the part still links.}

\engcheck{Customer view vs vendor view}{%
Vendor: internals, device physics, implementation.\\
Customer: BER, sensitivity, FEC, telemetry, environmental sweeps, interop.\\
Engineering samples (Tx-only, Rx-only, breakout, PRBS) open isolation;
otherwise stay on the external surface.}

For a bookended product, begin with BER, FEC, sensitivity, telemetry,
environment, and interop. Request engineering access only when that surface
cannot decide (\Cref{sec:tree-blackbox-access}).

Optoelectronics inherited a common qualification language from telecom:
\term{Telcordia GR-468-CORE}\firstcite{gr468}. The core stress tests still show up
on every laser and module program:

\begin{itemize}
\item HTOL (high-temperature operating life) for life or mechanism evidence.
\item Burn-in as a production infant-mortality screen when justified.
\item Temperature cycling and damp heat.
\item Electrostatic-discharge and mechanical stress.
\end{itemize}

Keep the jobs distinct: \textbf{burn-in} screens infant mortality from a
production population; \textbf{qualification HTOL} gathers life or mechanism
evidence under accelerated operation; a production burn-in is a manufacturing
screen only when separation, cycle time, and cost justify it. Do not imply that
every GR-468-style HTOL is a per-unit screen.

\term{Arrhenius} acceleration underpins life projection only when the named
failure mechanism is temperature-accelerated in the assumed regime: raising
temperature accelerates wear-out by a factor set by the activation energy for
that mechanism. Do not apply one $E_a$ to mixed mechanisms.

\paragraph{GR-468 in practice.}
\label{sec:gr468}

Telcordia GR-468-CORE is the common qualification language for optoelectronic
modules and discrete lasers. Map each stress onto the qualification evidence
path in \Cref{sec:tree-qual-evidence}; do not invent a second sequence here.

A 1{,}000-hour life test may justify a 100\% room-temperature proxy, a sampled
hot audit, a process monitor, or no direct production screen at all. Do not map
every GR-468 stress sequence into 100\% ATP. Document $E_a$ and confidence
bounds when converting HTOL hours to field years, keep sample-size humility
(\Cref{sec:fit-dppm,tab:npi}), and qualify the laser die, hermetic package, and
module assembly separately when failures split across those boundaries
(\Cref{sec:photonic-packaging,sec:laser-aging,sec:elsfp}).

\paragraph{Qualification planning matrix.}
\label{sec:qual-planning-matrix}

Qualification engineering starts from mechanisms, not from a museum of tests.
The matrix is illustrative: fill cells for the product class and claimed life.

\begin{table*}[ht]
\refstepcounter{table}\label{tab:qual-planning-matrix}%
\centering
\setlength{\tabcolsep}{3pt}%
\footnotesize
\renewcommand{\arraystretch}{1.15}%
\begin{tabular}{@{}p{0.18\linewidth}p{0.18\linewidth}p{0.18\linewidth}p{0.18\linewidth}p{0.20\linewidth}@{}}
\toprule
Failure mechanism & Stress & Observable & Acceptance & Production control \\
\midrule
Laser degradation & Temperature / lifetime (HTOL) & Power, wavelength, BER & Named limit vs life claim & ATP / SPC / burn-in proxy \\
\addlinespace[0.7em]
Solder fatigue & Temperature cycling & Resistance, BER, opens & Post-stress continuity / BER & Process control, FAIR \\
\addlinespace[0.7em]
Contamination / corrosion & Humidity / damp heat & Loss, ORL, leakage & IL/ORL / functional limits & Handling, sealing, audit \\
\addlinespace[0.7em]
Connector wear & Mate cycling & Insertion loss, ORL & Cycle-count IL budget & Supplier / hygiene control \\
\bottomrule
\end{tabular}
\par\medskip
{\raggedright\footnotesize\textbf{Table~\thetable.} Qualification planning
matrix. Each row is mechanism $\rightarrow$ stress $\rightarrow$ observable
$\rightarrow$ acceptance $\rightarrow$ production control
(\Cref{sec:tree-qual-evidence}; interview form \Cref{fw:qual-plan}).
Decision unlocked: which mechanism-stress-observable row is missing before you
claim life.\par}
\end{table*}

Interview and wall-chart form of the same path:
\Cref{fw:qual-plan,sec:tree-qual-evidence}.

\paragraph{Sample strategy and confidence.}
\label{sec:qual-sample-strategy}

Weak answer: ``We test 20 units.'' Strong answer: the sample strategy depends
on the failure-rate target, the confidence requirement, cost, and population
variation. State lots, date codes, suppliers, manufacturing sites or lines,
process corners, and whether the claim is zero-failure upper-bound or observed
rate (\Cref{sec:fit-dppm,sec:tree-evidence-block}). A rate without population,
observation time, and confidence is not a fleet claim.

\heuristic{A small sample that saw zero failures does not prove field life. It
sets an upper bound. State that bound, or do not make the claim.}

\tradeoff{More qualification vs faster release}{Confidence and fewer late
escapes}{Schedule, cost, and delayed learning in the field}{When remaining risk
is high and reversible controls cannot cover it}{Prioritize tests by risk
$\times$ uncertainty $\times$ impact. Do not test everything equally.}

\paragraph{GR-1221: the passive-component companion.}
GR-468 covers active optoelectronics. Its companion, \term{Telcordia GR-1221-CORE}
(Generic Reliability Assurance Requirements for Passive Optical Components), covers
the parts GR-468 does not: connectors, fiber couplers, WDM filters and MUX/DEMUX,
splitters, and isolators~\firstcite{gr1221}. It uses the same style of stress
sequence (damp heat, temperature cycling, mechanical, and aging tests) but scores
pass/fail on insertion loss and return loss rather than on LIV. A short-reach link
that leans on an on-package or blind-mate MUX and on external multi-wavelength
sources carries a passive reliability budget that lives in GR-1221, not GR-468
(\Cref{sec:connector-reliability,ch:wdm}). Split the qual the same way you split
the FIT: active laser die under GR-468, silicon under JESD47, passive optics under
GR-1221.

\paragraph{ATP sketch: EML module or ELSFP.}
A short acceptance sketch lives with the qual hooks in \Cref{sec:elsfp}; the full
ATP-as-contract, SPC, and 8D workflow is in \Cref{sec:supplier-exec}. Failures that
pass qual but fail field usually sit in derating policy or connector contamination
(\Cref{sec:laser-aging,sec:fleet-triage}).

\section{Electronics reliability: driver, TIA, and DSP silicon}
\label{sec:ic-reliability}

GR-468 covers the optoelectronic parts of the link: the laser die, the
photodiode, and the hermetic or non-hermetic package around them. The modulator
driver, TIA, retimer, and DSP (\Cref{sec:drivers,sec:pd-tia}) are ordinary CMOS or
SiGe BiCMOS ICs, and they wear out and fail by a different, better-documented set
of mechanisms. Treat them with the semiconductor industry's own qualification
language, not with Arrhenius laser-aging math borrowed from \Cref{sec:laser-aging}.

\paragraph{JESD47: the silicon-side GR-468.}
\firstterm{JESD47} JEDEC JESD47 is the baseline stress-test-driven qualification
flow for a new IC, a device family, or a process change: temperature cycling,
HTOL, HTSL (high-temperature storage life), autoclave or HAST (highly accelerated
stress test) for moisture, and mechanical shock and vibration~\firstcite{jesd47}.
It plays the same role for driver and TIA silicon that GR-468 plays for the
laser: a common list of stresses that a supplier runs once and a customer accepts
instead of renegotiating a qual plan on every design win.

\paragraph{ESD and latch-up: failure modes GR-468 does not test.}
Two mechanisms are specific to ICs and absent from the laser-side wear-out map in
\Cref{tab:wearout-map}:

\begin{description}
\item[ESD] a discharge event during handling or assembly damages a gate oxide or
      junction. Component-level classification uses the human-body model (HBM) and
      charged-device model (CDM) test standards, \term{ANSI/ESDA/JEDEC JS-001} and
      \term{JS-002}~\firstcite{jsesd}. A driver or TIA datasheet HBM/CDM rating is
      the number that protects the part on the factory floor, at fiber-attach and
      wire-bond stations where a laser die is also exposed.
\item[Latch-up] a parasitic thyristor structure in CMOS turns on under an
      overvoltage or current-injection event and holds a low-impedance path until
      power is cycled. \term{JESD78}\firstterm{JESD78} defines the overvoltage and
      $\pm100$~mA current-injection test that classifies susceptibility by supply
      and signal pin~\firstcite{jesd78}. A latched driver IC can look like a dead
      laser on the bench (no light, no LIV signature) until you check the supply
      current instead of the optical path.
\end{description}

Both mechanisms are 100\%-screen or design-margin items, not something you
project with an activation energy. If a driver fails ESD or latch-up in the
field, that is a manufacturability or design-margin bucket item
(\Cref{sec:fleet-triage}), not a wear-out FIT argument.

\paragraph{AEC-Q100: a borrowed grade, not a requirement.}
\term{AEC-Q100}\firstterm{AECQ100} is the automotive industry's qualification
standard for ICs, built on the same JEDEC JESD47/JESD22 stress methods with
tighter ESD targets and named temperature grades from Grade~3
($-40$ to $85$\textdegree C) up to Grade~0 ($-40$ to $150$\textdegree C)~\firstcite{aecq100}.
Datacenter optics does not require Q100; the fleet lives in a controlled data hall,
not an engine bay. It is still a useful signal: a driver, TIA, or retimer die that
also ships in an automotive part number typically carries a published Q100 grade,
and that grade is a fast proxy for the ESD/latch-up margin and temperature-cycle
depth behind the datasheet, without re-running the qual plan yourself.

\paragraph{Where this lands in the ATP.}
Fold IC-level qual into the same acceptance and SPC structure used for the laser
(\Cref{tab:atp-laser,sec:supplier-exec}): require the supplier's JESD47 qual
report and HBM/CDM/latch-up ratings for driver and TIA die at DVT, add an ESD
handling audit to the incoming-QC checklist alongside laser LIV/SMSR sampling,
and treat a driver/TIA silicon revision the same way you treat a laser die
revision or a CMIS firmware rev: an ECO that needs first-article requalification,
not a silent BOM swap.

\section{Wear-out modes to know}
\label{sec:wearout-modes}

Arrhenius math, derating, and the worked FIT example live in
\Cref{sec:laser-aging}. This section is the mechanism catalog: how each failure
shows up in ATP and telemetry, and which triage bucket owns it
(\Cref{sec:fleet-triage}). Do not run process CAPA on a wear-out part, and do not
burn FIT math on a dirty connector.

\paragraph{Infant mortality versus wear-out versus packaging.}
Field failures come in three clocks, and mixing them up wastes CAPA. Infant
mortality is early fails from latent defects; burn-in and HTOL screens remove them
before ship (\Cref{sec:gr468}). Wear-out is gradual or sudden end-of-life in the
laser or EAM under temperature, current, and optical power, projected with
Arrhenius and derating (\Cref{sec:laser-aging}). Packaging and assembly faults
(FAU align, epoxy creep, solder voids, connector wear) often dominate field
returns once lasers are screened (\Cref{sec:photonic-packaging}).
Destructive physical analysis (facet cross-section, EDX, FAU section) is required
when the signature is ambiguous or when you need evidence for supplier 8D
(\Cref{sec:supplier-exec}).

\paragraph{Mechanism map.}
\Cref{tab:wearout-map} is the working list for laser-bearing modules and CPO/ELS
paths. Customize limits in the ATP; keep the classification discipline.

\begin{table*}[ht]
\refstepcounter{table}\label{tab:wearout-map}%
\centering
\setlength{\tabcolsep}{3pt}%
\footnotesize
\renewcommand{\arraystretch}{1.25}%
\begin{tabular}{@{}p{0.16\linewidth}p{0.22\linewidth}p{0.28\linewidth}p{0.26\linewidth}@{}}
\toprule
Mechanism & Observable & ATP / telemetry & Triage bucket \\
\midrule
COD (facet) & Sudden dark or hard fail; was healthy & Dark LIV; DPA facet; date-code cluster? & Reliability (COD) or mfg (ESD) \\
\addlinespace[1.5em]
Gradual facet / active region & $I_\mathrm{th}$ up, slope down over life & LIV trend vs ship ATP; HTOL lot history & Reliability (wear-out) \\
\addlinespace[1.5em]
SMSR collapse & Side modes rise; modal noise / BER & OSA SMSR vs floor at $T$ & Reliability; watch aging \\
\addlinespace[1.5em]
EAM aging (EML) & TDECQ/RLM creep at fixed bias & EAM bias sweep + DCA; bias creep log & Reliability (EAM) \\
\addlinespace[1.5em]
RIN rise & BER floor up; feedback sensitive & RIN @ ORL; isolator / connector check & Perf if ORL; reliability if isolator \\
\addlinespace[1.5em]
TEC / thermal control & Unlock or $\lambda$ walk; LIV may look fine & TEC current, case $T$, lock status & Perf (lock) or reliability (TEC) \\
\addlinespace[1.5em]
Coupling / FAU / solder & Loss step, intermittent LOS, shock-related & ORL, mate cycles, DPA FAU/solder & Manufacturability / packaging \\
\addlinespace[1.5em]
Driver/TIA latch-up (ESD) & Sudden hard fail; no light, no LIV signature; supply current spikes & Supply current vs bias; JESD78 rating; date-code cluster? & Mfg (ESD) or design margin \\
\addlinespace[1.5em]
Connector wear / contamination & ORL creep after repeated mate cycles; RIN floor rise & Mate-cycle count vs IEC~61300-2-2 rating; endface grade (IEC~61300-3-35) & Manufacturability / packaging \\
\bottomrule
\end{tabular}
\par\medskip
{\raggedright\footnotesize\textbf{Table~\thetable.} Wear-out and packaging mechanisms
versus observables. Arrhenius projection and derating for the laser rows:
\Cref{sec:laser-aging}. Electronics stress qualification:
\Cref{sec:ic-reliability}. Connector reliability: \Cref{sec:connector-reliability}.
Field classification workflow: \Cref{sec:fleet-triage}.\par}
\end{table*}
\FloatBarrier

\subsection{Reading the wear-out map}
\label{sec:wearout-map-read}

\Cref{tab:wearout-map} is a triage map, not a life calculator. Classify before
corrective action. Reliability rows need life models and derate. Performance
rows need plant and control fixes. Manufacturing and packaging rows need process
and ATP screens. Mixing buckets wastes CAPA (\Cref{sec:fleet-triage}). Arrhenius
projection lives in \Cref{sec:laser-aging}; this section teaches how to read the
map in three families.

\paragraph{Semiconductor wear.}
COD, gradual active-region degradation, SMSR collapse, and EAM aging change the
diode or modulator physics. COD appears as sudden dark after a healthy ship LIV;
gradual wear appears as rising $I_\mathrm{th}$ and falling slope; SMSR collapse
raises modal noise while average power can still look fine; EAM aging creeps
TDECQ/RLM at fixed laser bias. Separate them with dark LIV and facet DPA (COD
versus ESD), LIV trends versus ship ATP (wear), OSA SMSR versus temperature and
age (modal), and EAM bias sweeps with eye metrics (modulator). Controls are life
review and derate for wear, SMSR/aging ATP for modal risk, and EAM life/ATP
screens for modulator aging. Do not burn Arrhenius math on an ESD lot, and do
not replace ``good'' DFBs while the EAM curve walks the eye closed.

\paragraph{Control and electronics.}
TEC/lock faults, driver/TIA latch-up or ESD, and ORL-driven RIN rise change the
environment or the electronics path more often than the laser die. Unlock or
$\lambda$ walk with healthy LIV points at cooler or lock; sudden hard fail with
supply-current spikes and no optical aging signature points at ESD/latch-up; a
BER floor that tracks ORL points at reflections before isolator death. Separate
them with TEC current, case $T$, and lock status; supply current and JESD
ratings with date-code clusters; RIN at stated ORL versus clean plant
(\Cref{sec:rin-values,sec:ic-reliability,sec:lock-validation}). Controls are
lock/thermal policy, ESD handling screens, and connector/reflection budget, not
laser FIT CAPA for every soft floor.

\paragraph{Packaging and optical interfaces.}
Coupling, FAU/solder, and connector wear or contamination produce step loss,
intermittent LOS, shock correlation, or ORL creep after mates. Semiconductor
metrics often stay clean. Separate them with ORL, mate-cycle history, endface
grade (IEC~61300), and DPA of FAU/solder when needed
(\Cref{sec:photonic-packaging,sec:connector-reliability}). Controls are assembly
CAPA, attach screens, service hygiene, and connector ratings. Life models must
not absorb packaging escapes.

Later evidence must not rewrite an earlier wrong bucket without new data. A
clean facet DPA does not clear a connector if ORL was never measured. An HTOL
pass does not clear a date-code ESD cluster.

\section{The reliability bathtub: three failure regimes}
\label{sec:bathtub}

Field failures follow three regimes with different clocks and different fixes.
Mixing them wastes corrective action on the wrong mechanism.

\begin{description}
\item[Infant mortality (early life)] Latent defects from manufacturing: weak
      solder joints, marginal laser die, contamination, firmware bugs that trip on
      first thermal cycle. Burns down rapidly with time. Fixed by burn-in,
      screening, tighter ATP, and first-article control (\Cref{sec:gr468}).
\item[Useful life (constant rate)] Random failures at a roughly steady FIT:
      cosmic-ray-induced single-event upsets, handling damage during unrelated
      service actions, and isolated material defects that passed all screens. Fixed
      by redundancy, field-replaceable modules, and design margin. Note: connector
      wear and contamination accumulate with mate cycles and are usage-driven
      degradation (\Cref{sec:connector-reliability}), not constant-rate random
      failures; budget them separately from the flat-hazard FIT.
\item[Wear-out (end of life)] Physics-driven degradation: laser facet, active
      region, EAM absorption curve shift, TEC aging, epoxy creep. Onset depends
      on temperature, current, optical power, and time. Fixed by derating,
      Arrhenius-based life projection, and planned replacement intervals
      (\Cref{sec:laser-aging}).
\end{description}

A rising failure rate after years of service is wear-out and calls for replacement
planning, not supplier 8D. A cluster of early failures on a new lot is infant
mortality and calls for tighter screens. A steady trickle with no date-code
correlation is useful-life random and calls for redundancy and fast repair. The
triage tree in \Cref{sec:fleet-triage} forces this classification before
corrective action starts.

\section{Yield analysis}
\label{sec:yield-analysis}

Yield is not one number. It splits by stage and by failure mode, and each split
points at a different owner.

\begin{description}
\item[Wafer / die yield] Process-limited: waveguide loss, ring resonance spread,
      heater shorts, photodiode dark current. Caught at wafer probe. Owner: foundry
      SPC.
\item[Assembly yield] Packaging-limited: fiber-array attach alignment, solder voids,
      wirebond pull, epoxy placement. Caught at module ATP. Owner: assembly supplier.
\item[Test yield (first-pass)] ATP-limited: units that fail one or more acceptance
      criteria on first pass. May include measurement-system false rejects (gauge
      RR). Owner: test engineering.
\item[Escaped DPPM (post-screen field failures)] Units that passed all screens but
      fail in the fleet. Each escape is either a preventable coverage gap (the
      screen exists but missed the defect) or a residual latent failure that no
      cost-effective screen separates from good units. Owner: quality and
      reliability engineering.
\end{description}

\begin{table*}[ht]
\refstepcounter{table}\label{tab:yield-split}%
\centering
\setlength{\tabcolsep}{3.5pt}%
\footnotesize
\renewcommand{\arraystretch}{1.25}%
\begin{tabular}{@{}p{0.22\linewidth}p{0.28\linewidth}p{0.22\linewidth}p{0.22\linewidth}@{}}
\toprule
Yield stage & Main limit & First catch & Owner \\
\midrule
Wafer / die & Waveguide, resonance, heater, PD dark & Wafer probe & Foundry SPC \\
\addlinespace[1.5em]
Assembly & FAU align, solder, wirebond, epoxy & Module ATP & Assembly supplier \\
\addlinespace[1.5em]
Test (first-pass) & ATP fails; may include false rejects & ATP station + gauge RR & Test engineering \\
\addlinespace[1.5em]
Escaped DPPM & Passed screens; failed in fleet & Field RMA / triage & Quality / reliability \\
\bottomrule
\end{tabular}
\par\medskip
{\raggedright\footnotesize\textbf{Table~\thetable.} Yield stages, first catch, and
owner. Split escapes further in \Cref{tab:escape-classes}.\par}
\end{table*}
\FloatBarrier

Track yield by ATP row, lot, supplier site, tester, and date code. A yield drop
that correlates with one tester is likely a measurement problem. A yield drop that
correlates with one supplier lot is likely a process problem. A yield drop with no
observed correlation requires further investigation: verify gauge repeatability,
expand stratification (shift, fixture, material lot, firmware), and test guardband
or specification mismatch as hypotheses before concluding. Do not open supplier
corrective action until the measurement system is cleared (\Cref{sec:hvm-test}).

\heuristic{Clear the tester with a golden unit before you escalate a supplier.
Station drift masquerades as a process excursion more often than engineers
admit.}

\usuallymeans{A golden unit fails on only one production station}{fixture,
calibration, cable, software limit, or operator path on that station}{a sudden
die-level failure of every good unit that station has ever seen}

\section{Escaped defect analysis}
\label{sec:escaped-defects}

\begin{dectree}
Supplier escape containment flow
  |
Escape detected
  |
Provisional containment
  |
Scope analysis
  |
Refine contained population
  |
Investigate / confirm mechanism
  |
Corrective action
  |
Recurrence control
\end{dectree}

Provisional containment, scope refinement, mechanism confirmation, and
recurrence control are different actions. Apply
\Cref{sec:tree-escape,sec:tree-recurrence,sec:tree-evidence-block} to separate
immediate hold from FA and from the ATP, sample, SPC, or telemetry change that
closes the loop. Organize the spent margin with the five ledgers before naming
a component (\Cref{sec:laser-margin-erosion,sec:debug-fork}).

An escaped defect is a unit that passed every production screen and failed in the
field. Post-screen field failures split into two categories with different
corrective actions:

\begin{table*}[ht]
\refstepcounter{table}\label{tab:escape-classes}%
\centering
\setlength{\tabcolsep}{3.5pt}%
\footnotesize
\renewcommand{\arraystretch}{1.25}%
\begin{tabular}{@{}p{0.20\linewidth}p{0.28\linewidth}p{0.24\linewidth}p{0.22\linewidth}@{}}
\toprule
Class & Meaning & Typical action & Lands in \\
\midrule
Preventable coverage & Screen could have caught it & Add/tighten ATP or SPC & Escape DPPM, CAPA \\
\addlinespace[1.5em]
Residual latent & No cost-effective screen & FIT / redundancy / replace & Residual FIT model \\
\bottomrule
\end{tabular}
\par\medskip
{\raggedright\footnotesize\textbf{Table~\thetable.} Escape classes. Preventable rows
change production; residual rows change the life model.\par}
\end{table*}
\FloatBarrier

\paragraph{Preventable coverage escapes.}
A defect that a cost-effective screen could have caught but did not, because the
screen was absent, miscalibrated, or insufficiently stressed. Each one requires an
ATP or process-control change. Common mechanisms in optical modules:

\begin{itemize}
\item \textbf{Contamination:} particles trapped during assembly that only move under
      thermal cycling or vibration, shifting coupling or raising ORL intermittently.
\item \textbf{Marginal alignment:} fiber-attach or FAU position within spec at room
      temperature but outside at the thermal corner, creating a
      temperature-dependent loss that the ATP chamber did not run.
\item \textbf{Weak solder or wirebond:} passes pull test but fatigues under
      repeated thermal cycling, creating intermittent opens or increased resistance.
\item \textbf{Firmware corner:} a state-machine path or calibration table entry that
      production never exercises but the fleet hits on a specific boot sequence,
      temperature bin, or host interaction.
\item \textbf{Packaging stress:} residual mechanical stress from underfill cure or
      lid attach that relaxes over time, shifting alignment or birefringence.
\end{itemize}

For each preventable escape, trace the failure signature back to the earliest point
in the production flow where it could have been caught, and add or tighten the
screen there.

\heuristic{An escape without an ATP or SPC change is unfinished work. Containment
stops the bleed; the control stops the next lot.}

\paragraph{Residual latent failures.}
A defect that no cost-effective screen can separate from good units at the time of
test. Examples include cosmic-ray-induced single-event upsets, rare material
inclusions below detection limits, and process-marginal units that sit inside all
guardbands but fail under a specific field combination of temperature, ORL, and
neighbor load. These go into the residual FIT model, not into ATP. Document why no
reasonable screen exists so the decision is auditable and revisitable as test
technology improves.

\section{Photonic packaging and module-level failures}
\label{sec:photonic-packaging}

Fleet FIT is not only laser wear-out. Once lasers are screened and derated,
module and packaging failures often dominate field returns: the part that shipped
with a clean LIV still loses light after shock, humidity, or a thousand ELSFP mate
cycles. Fiber attach and FAU alignment fail from shock, humidity ingress, and
epoxy creep; CPO fiber-array units add assembly steps that wafer test cannot catch
(\Cref{sec:cpo-status}). Hybrid stacks (TFLN-on-Si, InP laser on Si, flip-chip
drivers) introduce solder voids, underfill cracks, and RF return-loss drift
(\Cref{sec:tfln-mzm,sec:drivers}). Thermal paths matter too: uncooled datacom
versus liquid-cooled XPO/CPO, and TEC failure that looks like wavelength drift off
grid or off ring (\Cref{sec:siring,ch:wdm}).

\paragraph{Connector reliability: MPO, mating cycles, and endface quality.}
\label{sec:connector-reliability}

Multi-fiber connectors are the highest-touch mechanical interface in the fleet:
every ELSFP swap, every fiber-attach unit (FAU)\firstterm{FAU} rework, and every cable-plant
install mates and unmates an MPO\firstterm{MPO}. The MPO/MT ferrule family (rectangular, 6.4~mm
$\times$ 2.5~mm, guide-pin aligned, 8/12/16/24 fibers per row) is standardized in
\term{IEC 61754-7}, split into one-fibre-row and two-fibre-row
parts~\firstcite{iec617547}. That standard fixes geometry, not lifetime; lifetime
comes from two companion test methods. \term{IEC 61300-2-2} specifies the
mate/unmate cycling test connector datasheets are rated against, and \term{IEC
61300-3-35} grades endface scratches, pits, and debris into pass/fail zones on the
fiber core and cladding~\firstcite{iec61300}. TIA-568.3 sets 500 cycles as the
structured-cabling
mating-durability floor; MPO/MTP-class connectors in practice are commonly rated
well above 1000 cycles, but that headroom erodes fast with the wrong cleaning
discipline (\Cref{sec:optical-channel}).

Three practical consequences follow for an ELSFP or CPO fiber-attach program.
First, ORL creep is a mating-cycle and cleaning problem before it is a laser
problem: a rising RIN floor after repeated ELS swaps (\Cref{tab:wearout-map}) is
diagnosed with an IEC~61300-3-35-style endface inspection, not a laser FA request.
Second, mate-cycle count belongs in the same telemetry you already read for CMIS
and DDM (\Cref{sec:cmis}); track it per connector, not per module, since a
connector can outlive several module swaps or vice versa. Third, write the
mating-cycle and endface-grade limits into the ATP explicitly
(\Cref{tab:atp-laser}) rather than inheriting a generic MPO datasheet number: an
ELS bank that hot-swaps weekly reaches a 500-cycle floor in under ten years, and a
CPO fiber array that is field-serviced more aggressively reaches it faster still.

ELSFP cycling adds connector wear and contamination that raise ORL
(\Cref{sec:optical-channel,sec:elsfp}); the mating-cycle and endface-grade limits
above are exactly the numbers that turn ``the connector feels loose'' into an ATP
line item instead of a guess.

Destructive physical analysis (cross-section, EDX) and structured 8D/CAPA with
suppliers close the loop from RMA to design rule
(\Cref{sec:supplier-exec,sec:fleet-triage}). Without that loop, packaging FIT gets
mis-attributed to laser Arrhenius models and the wrong part gets redesigned.

\section{Production test at volume}
\label{sec:hvm-test}

\engcheck{Before production}{%
ATP $\cdot$
SPC $\cdot$
telemetry $\cdot$
supplier gates $\cdot$
monitoring owners $\cdot$
RMA-to-ATP feedback
(\Cref{sec:tree-checklists}).}

\subsection{Test time is a cost, coverage is a risk}

Every second in the acceptance test plan (ATP) times millions of units is line
capacity and real money. Every skipped measurement creates uncontrolled escape
risk; it is not automatically a field DPPM event (\Cref{sec:fit-dppm}). The core
tension in high-volume manufacturing is how much coverage you buy per second. The
expensive optical steps are thermal soak and corner runs, TDECQ on a sampling
scope, BER dwell long enough to trust a low pre-FEC target, laser burn-in, and
mate-cycle stress on ELSFP connectors. Some screens are statistical (sample
burn-in from a lot, audit TDECQ on a subset). Safety and enable-sequence faults
usually require 100\% coverage. For source-level production where LIV or SMSR
are directly available and correlated to escape risk, they may be 100\% screens.
For closed modules, use the validated module-level proxy or a documented sampling
plan (\Cref{sec:cmis,sec:laser-safety,tab:atp-laser}).

\tradeoff{More production screening vs cost}{Escape detection and earlier
catch}{Cycle time, tester cost, and false rejects that burn good units}{When a
named mechanism has a cheap, reliable detection signature}{Choose the cheapest
control that reliably detects the failure mode: 100\% ATP, sample audit, SPC,
or supplier process control.}

\tradeoff{Burn-in vs cycle time}{Infant-mortality removal when the screen
separates}{Line capacity, cost, and stress on healthy units}{When escape data
and mechanism justify the screen on this population}{Keep burn-in only while it
buys escapes you cannot catch cheaper elsewhere.}

\subsection{Where the test happens: wafer, die, module, system}

Push defect detection as far upstream as correlation allows. Wafer-level or PIC
probe catches process shifts (waveguide loss, ring resonance drift, bad heaters)
before fiber attach and packaging spend. Killing a bad die at probe is orders of
magnitude cheaper than an RMA (\Cref{sec:photonic-packaging}). Module ATP is the
full functional test: optical power class, TDECQ or proxy, sensitivity spot-check,
CMIS bring-up, and connector/ORL on ELS parts. System or golden-host bring-up
catches interop: media type, firmware rev, equalizer defaults, and the corners in
\Cref{sec:prod-corners}. Wafer test cannot catch fiber attach, FAU alignment,
epoxy creep, or connector wear. Those failures must survive to module ATP and,
for some signatures, to fleet telemetry (\Cref{sec:fleet-triage}).

\subsection{ATE-to-bench correlation and gauge R\&R}
\label{sec:gauge-rr}

Production testers are built for speed and cost, not lab fidelity. Correlation
asks whether ATE TDECQ, OMA, or sensitivity tracks the DCA/BERT reference within
a known offset and spread. Teach the measurement system explicitly:

\begin{description}
\item[Repeatability] Same station, operator, and unit.
\item[Reproducibility] Across stations, shifts, and operators.
\item[Bias] Production station versus reference lab.
\item[False reject / false accept] Guardband from measurement uncertainty.
\item[Golden-unit stability and station drift] Catch a stale golden or drifting
      ATE before it becomes a yield cliff or DPPM escape.
\end{description}

Keep a golden module (and golden laser subassembly for ELS), run gauge R\&R
across testers and shifts, and correlate CMIS monitors to bench instruments the
same way you correlate TDECQ (\Cref{sec:cmis}). If the ATE and the DCA disagree,
fix the correlation before you argue with the supplier about spec.

\subsection{Screens, guardbanding, and SPC}

Burn-in (infant-mortality screen) and HTOL (life/mechanism evidence) trade
different risks against test time and cost
(\Cref{sec:wearout-modes,sec:gr468}). Test limits are usually guardbanded
tighter than the customer spec so field DPPM stays inside target under drift.
SPC control charts on LIV, SMSR, RIN, TDECQ, and mate-cycle yield by lot, site,
and date code catch a process shift before it becomes an 8D (\Cref{sec:supplier-exec}).
Production test is a yield, DPPM, and cost trade under a fixed reliability target.
It is not a pass/fail checkbox after the optics already work on a golden bench.

\section{Supplier execution playbook}
\label{sec:supplier-exec}

The supplier path is milestones, performance targets, quality, and
manufacturability triage. That is not a soft skill. It is a concrete contract:
requirements, gates, acceptance tests, process control, and corrective action when
a lot goes wrong.

\whyexp{care about production lots?}{Because manufacturing escapes almost always
correlate with process history. Lot, date code, site, and firmware tags often
beat another night on one returned unit.}

\heuristic{Ask for the process change list before you invent new physics. Most
lot escapes sit next to a real change record.}

\tradeoff{Second source vs qualification burden}{Supply resilience and pricing
leverage}{Validation, interop matrix, and manufacturing differences}{When supply
or concentration risk exceeds the qual cost}{Qualify second sources based on
risk and evidence, not ideology.}

Lab hero samples versus manufacturable yield: same tradeoff as in
\Cref{ch:validation}. Optimize the system you can build, not the best bench unit.

\begin{dectree}
Design requirements
  |
Qualification
  |
ATP
  |
Production data
  |
Fleet data
  |
Failure analysis
  |
Updated limits or screens
  |
Next production cycle
\end{dectree}

Production validation is replayable and decision-oriented
(\Cref{sec:tree-production-loop}).

\paragraph{NPI gates and exit criteria.}
New product introduction (\term{NPI}\firstterm{NPI}) gates are the manufacturing
face of the validation ladder. Write exit criteria a supplier can fail without
ambiguity, not slogans. EVT/DVT/PVT/MP are stage names, not a calendar; dates and
sample sizes belong in the program plan.

\paragraph{Why program maturity changes the question.}
EVT asks whether the architecture can be made to work. DVT asks whether its
behavior and margin are understood across intended corners. Qualification asks
whether named failure mechanisms threaten life. PVT asks whether production
tooling and suppliers reproduce the qualified result. A pilot asks whether
laboratory and factory assumptions survive deployment. MP is not another
validation experiment; it is sustained control with SPC, ECO, and RMA ownership
(\Cref{sec:ladder-mp,sec:ladder-fleet-monitor}).

Do not use an EVT hero sample as PVT evidence, and do not treat MP as fleet
monitoring alone. Hold a gate if the exit data are missing
(\Cref{tab:npi,sec:validation-workflow}).

\begin{table*}[ht]
\refstepcounter{table}\label{tab:npi}%
\centering
\setlength{\tabcolsep}{4pt}%
\footnotesize
\renewcommand{\arraystretch}{1.25}%
\begin{tabular}{@{}p{0.10\linewidth}p{0.22\linewidth}p{0.36\linewidth}p{0.26\linewidth}@{}}
\toprule
Gate & Question & Representative evidence & Release decision \\
\midrule
EVT & Does it operate at all? & First light; CMIS bring-up; basic LIV/SMSR/RIN; one link closes BER & Continue / redesign integration \\
\addlinespace[1.25em]
DVT & Does it meet spec across corners? & Full ATP at $T$/$V$; prod-rep corners; stress plan + FIT model frozen & Enter qual / PVT / hold \\
\addlinespace[1.25em]
Qual & Env / reliability evidence ready? & Named mechanisms; sample plan; confidence (\Cref{sec:tree-qual-evidence}) & Enter PVT / hold \\
\addlinespace[1.25em]
PVT & Is it buildable at yield? & Multi-lot yield; SPC; burn-in escape; FAIR; production-host bring-up & Enter pilot / hold \\
\addlinespace[1.25em]
Pilot & Do assumptions hold in a bounded field trial? & Known serials/lots; enhanced telemetry; exit/rollback criteria & Open MP / restrict \\
\addlinespace[1.25em]
MP & Is quality sustained? & Steady DPPM; owned RMA Pareto; ECO control; fleet feedback & Keep shipping / CAPA / restrict \\
\bottomrule
\end{tabular}
\par\medskip
{\raggedright\footnotesize\textbf{Table~\thetable.} NPI gates. Decision unlocked:
which release call the gate evidence supports. Pilot sits between PVT and MP;
MP is sustained control, not fleet monitoring alone.\par}
\end{table*}

\paragraph{Requirements and ATP are the contract.}
ATP\firstterm{ATP} and the requirements doc are the contract. Write both and keep
them versioned together:

\heuristic{If a requirement has no ATP or sample line, it is a wish. If an ATP
line has no requirement, it is cost without a decision.}

\begin{enumerate}
\item \textbf{Requirements / PRD slice for the laser path:} fill
      \Cref{tab:laser-prd,sec:laser-reqs} (power class, grid, RIN@ORL, SMSR,
      derating, CMIS, FIT). Version it with the ATP.
\item \textbf{Acceptance test plan (ATP):} the measurable tests that prove those
      requirements on every ship lot (or on a defined sample). Map each ATP line to
      a GR-468 or design-validation stress where life is claimed (\Cref{sec:gr468}).
\end{enumerate}

\subsection{How production evidence is purchased}
\label{sec:atp-laser-read}

Every second of ATP costs line capacity. Every omitted screen leaves risk. Buy
evidence in categories that cannot substitute for each other.
\Cref{tab:atp-laser} is a working checklist for an EML pluggable or an ELSFP CW
module. Customize limits from the datasheet and the link budget; do not invent
numbers in the ATP itself.

\paragraph{Every-unit fast screens.}
Power class, CMIS state and bring-up, basic lane operation, and selected
electrical or optical checks that are cheap enough for 100\% coverage. These
catch dead units and firmware fails. They do not establish life.

\paragraph{Sampled expensive measurements.}
TDECQ, intrinsic and stressed RIN, detailed spectrum/SMSR, thermal corners, and
longer BER dwell. Use lot sample or audit when the signature is expensive and
correlated to escape risk.

\paragraph{Qualification-only stresses.}
HTOL, humidity, extensive cycling, and destructive analysis gather life or
mechanism evidence. They are not per-unit ship screens
(\Cref{sec:gr468,sec:tree-qual-evidence}).

\paragraph{Process controls.}
SPC on LIV, SMSR, RIN, TDECQ, and mate-cycle yield; golden-unit tracking; gauge
R\&R; incoming quality; first-article / FAIR after tooling or site change
(\Cref{sec:gauge-rr}).

\paragraph{Fleet controls.}
Telemetry, lot traceability, RMA codes split by mechanism, and recurrence
monitoring (\Cref{sec:fleet-triage}). These catch what ATP never saw.

A fast production power check does not establish life. HTOL does not prove every
shipped unit has correct firmware. SPC does not detect an unmeasured mechanism.
LIV, SMSR, and wavelength protect semiconductor identity; RIN and ORL protect
the reflection environment; EAM/DCA protects Tx quality on EML paths; CMIS
protects field evidence; burn-in and ESD protect infant mortality and handling;
thermal class protects the derate claim.

\begin{table*}[ht]
\refstepcounter{table}\label{tab:atp-laser}%
\centering
\setlength{\tabcolsep}{2.5pt}%
\footnotesize
\renewcommand{\arraystretch}{1.2}%
\begin{tabular}{@{}p{0.18\linewidth}p{0.16\linewidth}p{0.14\linewidth}p{0.22\linewidth}p{0.22\linewidth}@{}}
\toprule
Item & Method & Control class & Pass intent & Ties to \\
\midrule
LIV ($I_\mathrm{th}$, slope, kink) & SMU + power meter & 100\% ATP (source) / module proxy & kink-free bias window & wear-out, derate \\
\addlinespace[1.1em]
SMSR & OSA & 100\% ATP (source) / lot sample & single-mode vs.\ floor & modal noise \\
\addlinespace[1.1em]
Intrinsic RIN & PD + ESA & Lot sample / FA & quiet source floor & BER floor budget \\
\addlinespace[1.1em]
Stressed $\mathrm{RIN}_x\mathrm{OMA}$ & PD + ESA @ ORL & Lot sample / 100\% if escape & stressed Tx metric & named PMD \\
\addlinespace[1.1em]
Wavelength / grid & OSA / wavemeter & 100\% ATP or sample & channel ID & WDM lock \\
\addlinespace[1.1em]
Optical power class & power meter & 100\% ATP & class met & link budget \\
\addlinespace[1.1em]
EAM / TDECQ (EML) & bias + DCA & Lot sample / audit & ER, RLM, TDECQ & Tx quality \\
\addlinespace[1.1em]
CMIS / TWI bring-up & host / CMIS tool & 100\% ATP & state machine & telemetry \\
\addlinespace[1.1em]
Connector / ORL & mate + ORL meter & Periodic audit / sample & cycles + endface & packaging \\
\addlinespace[1.1em]
Burn-in (infant) & production screen & 100\% or lot sample & infant culled & not HTOL life \\
\addlinespace[1.1em]
HTOL life evidence & accelerated life & Qualification only & mechanism + FIT claim & GR-468 \\
\addlinespace[1.1em]
Driver/TIA ESD & JESD47 report & Qualification / audit & rating on file & IC reliability \\
\addlinespace[1.1em]
Thermal class & chamber & Lot sample / SPC & LIV/RIN/CMIS pass & derate \\
\addlinespace[1.1em]
Process monitors & SPC charts & SPC / process & drift detection & escape prevention \\
\addlinespace[1.1em]
Fleet cohort metrics & telemetry & Fleet telemetry & trend / alarm & recurrence \\
\bottomrule
\end{tabular}
\par\medskip
{\raggedright\footnotesize\textbf{Table~\thetable.} Control checklist for laser-bearing
modules (EML or ELSFP). Control class separates 100\% ATP, lot sample, audit,
qualification-only, SPC, and fleet telemetry. Intrinsic RIN and stressed
$\mathrm{RIN}_x\mathrm{OMA}$ are different metrics.\par}
\end{table*}
\FloatBarrier

\textbf{Exit for the ATP as a whole:} every ship lot (or defined sample) has
traceable pass data against versioned limits tied to the requirements slice.
\textbf{Decision:} ship, hold, or reopen DVT limits.

\paragraph{Incoming QC and SPC.}
Qual lots are small. Production catches drift that qual missed.

\begin{itemize}
\item \textbf{Incoming:} sample or 100\% screen against a subset of the ATP (at least
      power, CMIS, and a laser LIV/SMSR sample). Track DPPM by date code and site.
\item \textbf{SPC}\firstterm{SPC}: control charts on $I_\mathrm{th}$, slope, SMSR, Tx
      power, and burn-in fallout. A process shift is a hold, not a hope.
\item \textbf{First-article / FAIR:} when tooling, epi, or assembly site changes,
      rerun a defined FAIR package before open PO volume. Treat CMIS firmware revs
      the same way as process changes.
\end{itemize}

\paragraph{Excursions: 8D / CAPA.}
When a lot fails ATP, incoming, or field triage lands in the manufacturability
bucket (\Cref{sec:fleet-triage}), run structured corrective action:

\begin{enumerate}
\item \textbf{Contain:} quarantine WIP and ship holds; identify suspect date codes
      in the fleet.
\item \textbf{Evidence pack:} failing ATP rows, CMIS dumps, LIV/SMSR/RIN plots, DPA
      photos (facet, solder, FAU cross-section) compared to a golden unit.
\item \textbf{8D / CAPA}\firstterm{eightD}: confirmed mechanism with the supplier (process step,
      material lot, firmware), corrective action, and preventive control (ATP tighten,
      SPC limit, poke-yoke).
\item \textbf{Verify closure:} containment confirmed effective; mechanism reproduced
      or physically confirmed; corrective action removes the failure; no
      unacceptable regression introduced; production control detects recurrence;
      next lots remain stable; field cohort trend improves. Re-run FAIR alone is
      not enough for environmental, intermittent, or fleet-specific escapes
      (\Cref{sec:tree-evidence-block}).
\end{enumerate}

Do not close 8D on ``operator error'' without a control that would have caught it
at ATP or in process. If FA shows laser wear-out on a young unit, it may be a
reliability screen gap, not a supplier process bug; reclassify with
\Cref{sec:fleet-triage} before you argue FIT.

\paragraph{Milestone hygiene with partners.}
Align the partner calendar to gates, not slideware:

\begin{itemize}
\item freeze requirements before DVT samples are built;
\item freeze ATP limits before PVT yield is claimed;
\item freeze FIT/$E_a$ assumptions before reliability marketing numbers ship;
\item require ECO notice on laser die revision, TEC vendor, FAU epoxy, driver/TIA
      silicon revision (\Cref{sec:ic-reliability}), and CMIS firmware.
\end{itemize}

Your job in those meetings is to name the measurement that would kill the gate.
If nobody can point to an ATP row or a corner, the milestone is not real.

\keyidea{Reliability at scale is mechanism discipline plus supplier gates.
Classify with the wear-out map (\Cref{tab:wearout-map}) before FIT or 8D. Laser
wear-out uses Arrhenius and GR-468; driver/TIA uses JESD47 and ESD ratings
(\Cref{sec:ic-reliability}); connectors use mate-cycle and endface limits
(\Cref{sec:connector-reliability}). Gate suppliers on ATP, multi-lot SPC, and
FAIR. Do not run process CAPA on wear-out, or FIT math on a dirty connector.}

\section{From component FIT to fabric availability}
\label{sec:fabric-availability}

The FIT arithmetic in \Cref{sec:fit-example} gives a rate: about $0.6$ laser
failures per day for a fleet of $5\times10^5$ lasers at 50~FIT. That number sizes
the RMA pipeline and the ELS spares bin (\Cref{sec:elsfp}), but it does not say
what a failure costs or how a running job survives one. Two facts turn a
per-component rate into a fabric problem.

First, a training or large inference job is synchronous. A collective
(\Cref{sec:collectives}) waits for its slowest member, so a single dead or slow
link stalls the whole group, not just one endpoint (\Cref{sec:inference-bottlenecks}).
A link that flaps for a second is a stall for every accelerator in that collective.
The optical FIT the earlier chapters budget therefore matters out of proportion to
its share of the parts count.

Second, at cluster scale failures are continuous, not rare. Meta's published
Llama~3 run is the clearest public data point: 16{,}384 H100 GPUs over 54 days
logged 466 interruptions (419 unexpected), roughly one every three hours, while
holding about 90\% effective training time~\firstcite{metallama3}. GPU and HBM3
faults dominated at close to half; network switch and cable faults were 35 events,
8.4\% of the total. The optical link is a minority of hard job stops, but 8.4\% of
a failure every three hours is still tens of network events per run, and the ELS,
module, and connector FIT this chapter budgets
(\Cref{sec:fit-dppm,sec:connector-reliability}) lands in exactly that bucket.

So the design question shifts. It is no longer ``how reliable is one link'' but
``how does a fabric of $10^5$ links keep a job running through a failure every few
hours.'' The answers are architectural, and the optical engineer feeds each one.

\begin{description}
\item[Redundancy and rails.] Rail-optimized topologies (\Cref{sec:topologies})
      already give parallel planes; dual-plane and dual-ToR designs let a lost link
      degrade bandwidth instead of dropping an endpoint. Redundancy multiplies the
      link and laser count, which feeds straight back into the FIT budget: more
      resilience is more parts that can fail.
\item[Detection and reroute.] Transient faults stay below the job. KP4 FEC
      (\Cref{sec:kp4}) absorbs the error bursts a marginal link throws; link-level
      retry and sub-second link-flap detection plus adaptive routing steer traffic
      off a degraded link before the scheduler notices. Vendor fabrics (NVIDIA
      Spectrum-X and Quantum, Broadcom Tomahawk) advertise adaptive or ``cognitive''
      routing and link-level retry for this. Treat the specifics as vendor
      orientation, but the mechanism is why transient optical faults rarely reach
      the hard-stop bucket above.
\item[Topology reconfiguration.] When a link or rack dies for good, an optical
      circuit switch re-wires the topology around it in milliseconds, so the
      scheduler routes around the dead node instead of stalling the pod
      (\Cref{sec:ocs})~\citep{tpuv4ocs}. Component FIT still applies; the fabric
      survives each failure by re-wiring optically.
\item[Sparing and field service.] Hot spare nodes and lanes cover the interval
      between failure and repair. Field-replaceable external lasers (\Cref{sec:elsfp})
      make a dead laser a faceplate swap rather than a fabric outage, which is the
      architectural reason ELS decouples laser FIT from switch FIT. The connector
      mating-cycle and endface budget (\Cref{sec:connector-reliability}) sets how
      many of those swaps the plant survives.
\end{description}

The cost of a failure closes the loop. A hard interruption is lost compute plus
the time to detect, reroute or reschedule, and restart from the last checkpoint.
Fast detection and reroute shrink that lost time, which is the fabric-level reason
the module work in this chapter pays off: derating (\Cref{sec:laser-aging}),
burn-in and screens (\Cref{sec:gr468}), and a tight ATP (\Cref{sec:supplier-exec})
lower the failure rate, and a resilient fabric lowers the cost of each failure that
slips through. The two multiply.

\section{Engineering lens}
\label{sec:reliability-engineering-lens}

\subsection{How it works}

At fleet scale, reliability and manufacturability are design constraints: a modest
per-part FIT times millions of parts is a steady stream of failures. The chapter's
qualification, yield, and supplier discipline all aim at keeping that stream small
and classifiable.

\subsection{How it is measured}
\label{sec:reliability-how-measured}

Reliability is measured as a distribution over stress and time. Qualification
records failures by mechanism, stress, lot, and sample history. Production records
yield, fallout by ATP row, measurement distributions, gauge repeatability, and
escaped defects per million. Fleet data add install age, temperature, firmware,
supplier lot, return code, and no-fault-found rate. Keep the chain from wafer or
die measurement through module ATP to field return so a drift can be traced to its
first observable point (\Cref{sec:hvm-test,sec:supplier-exec,sec:fleet-triage}).

\subsection{How it fails}
\label{sec:reliability-how-fails}

Programs fail at scale through wear, variation, and escapes. Wear shifts a unit
after ship. Variation produces a weak tail across wafer, lot, site, or assembly
line. An escape is a defect the current screen cannot see or a control that was not
run. Yield can drop because the product changed, the process moved, incoming
material moved, calibration changed, or the tester moved. Do not open supplier
corrective action until the measurement system is cleared.

\failuremode{Yield drop}{First-pass yield falls from its stable baseline, often on
one ATP row, tester, lot, or shift.}{Process drift, supplier-lot variation,
assembly change, calibration or fixture drift, software revision, or a changed
guardband.}{Pareto by test and lot, golden-unit history, gauge repeatability,
station correlation, incoming data, and destructive analysis on selected
failures.}{Contain suspect material, clear the tester, identify the first changed
input, correct it, and verify with a controlled lot before release.}

\subsection{How it is debugged}
\label{sec:reliability-how-debugged}

For a yield fall, freeze software, limits, and suspect material. Split by tester,
shift, lot, supplier site, and ATP row. Golden unit across stations; failing unit
on a reference bench. Station-follows means repair the measurement system;
unit-follows means upstream process and mechanism FA. Contain first, confirm
second, change the process third, verify on fresh data.

\debugstory{Module yield fell sharply after a supplier lot change.}{The failure
Pareto pointed to one-lane TDECQ. Golden units passed all stations, and failed
units kept the bad lane on the reference bench. Cross-sections showed a shifted
fiber-array attach.}{The electrical path and testers were stable.}{An assembly
fixture change moved one fiber row outside its coupling window.}{The lot was held,
the fixture was restored, first-article coupling checks were tightened, and the
supplier control plan was revised.}

\section{Interview takeaway}
\label{sec:reliability-design-review}

\keyidea{Production readiness joins mechanism-based qualification, a stable
measurement system, controlled supplier changes, and field data that preserve lot
history. When yield or fleet health moves, contain the population, clear the test
system, find the first changed input, and verify the corrective control with new
data.}

Junior mistake: treat zero fails in a small HTOL lot as a FIT claim, or escalate
a supplier before the measurement system is cleared
(\Cref{sec:qual-sample-strategy,ch:failure-modes,app:case-studies}).

\paragraph{Three questions to test yourself.}
\begin{enumerate}
\item Why does qualification on a small sample not guarantee field reliability?
\item Yield dropped from 95\% to 70\% on one ATP row. What are your first three
      actions?
\item Which corrective action changes a control so the same defect cannot escape
      again?
\end{enumerate}

